\documentclass[12pt]{article}
\usepackage[utf8]{inputenc}
\usepackage[T1]{fontenc}
\usepackage{hyperref}
\title{Manado Malay Language Resources for K-12 Curriculum Development}
\author{Compiled by Claude Sonnet 5.5 (claude-sonnet-5-5) for Andr\'as Kornai}
\date{}
\begin{document}
\maketitle

\section{Introduction}
Manado Malay (Bahasa Manado; ISO 639-3 xmm; Glottocode mala1481) is a Malay-based creole spoken in Manado and the surrounding area of North Sulawesi, Indonesia. According to the Wikipedia article, it is used mainly in speech and has no standard orthography. Typical features include \emph{ba-} and \emph{ma(ng)-} in place of Indonesian \emph{ber-} and \emph{me(ng)-}, and the pronouns \emph{kita} (I), \emph{ngana} (you, singular) and \emph{ngoni} (you, plural).

\section{Primary Educational Resources}
No primers, textbooks or school programmes in Manado Malay turned up in the searches.

\section{Linguistic and Reference Materials}
\begin{itemize}
\item Interactive Manado dictionary on Webonary, with an introduction page: \url{https://www.webonary.org/manado/overview/introduction/}
\item Page on the history of the Manado language on the same site: \url{https://www.webonary.org/manado/?p=146806}
\item Ruben Stoel (2005), \emph{Focus in Manado Malay: grammar, particles and intonation}, the main book-length grammar (Glottolog reference entry): \url{https://glottolog.org/resource/reference/id/122282}
\item Overview article: \url{https://en.wikipedia.org/wiki/Manado_Malay}
\end{itemize}

\section{Religious and Community Resources}
None found in the searches.

\section{Digital and Audio Resources}
The Webonary dictionary is the only digital resource found. No audio or video collections turned up.

\section{Curriculum Development Context}
The language is not known to be taught in schools. Spelling is unstandardised, so a curriculum would first need an agreed orthography; the Webonary dictionary project supports development of a written form. Stoel's grammar and the dictionary are the base for content development. No community organisations were identified in these searches.

\end{document}